\chapter{Notation and Glossary}\label{app:notation}

This appendix collects all mathematical notation used in the monograph in a
single reference table, followed by a glossary of key technical terms.  Symbol
conventions are drawn directly from the paper \cite{kbound2026repo} and are
consistent throughout all chapters.

% -----------------------------------------------------------------------
\section{Notation table}
\label{sec:app-notation-table}
% -----------------------------------------------------------------------

\begin{longtable}{lp{11cm}}
\caption{Mathematical notation used throughout the monograph.}
\label{tab:notation}\\
\toprule
Symbol & Meaning \\
\midrule
\endfirsthead
\toprule
Symbol & Meaning \\
\midrule
\endhead
\midrule
\multicolumn{2}{r}{\textit{continued on next page}}\\
\endfoot
\bottomrule
\endlastfoot

$P_S(X,Y)$ &
  Source distribution.  Labels $Y$ are available at training and validation
  time. \\

$P_T(X,Y)$ &
  Target distribution.  At deployment only $P_T(X)$ is observable (no target
  labels). \\

$f_0$ &
  Frozen (baseline) predictor.  Deployed without modification unless
  \textsc{adapt} fires. \\

$f_a$ &
  Adapted (candidate) predictor.  The output of an adaptation mechanism
  (Tent, EATA, SAR, or any other) that is being evaluated. \\

$\ell$ &
  Per-sample loss function, e.g.\ $0/1$ loss or cross-entropy. \\

$R_T(f)$ &
  Target risk of predictor $f$: $R_T(f) = \E_{P_T}[\ell(f(X),Y)]$. \\

$\Delta$ &
  Adaptation benefit: $\Delta = R_T(f_0) - R_T(f_a)$.  Positive means
  adapting reduces risk (helps); negative means it hurts. \\

$\Delta(z)$ &
  Conditional benefit at evidence value $z$: $\Delta(z) = \E[\ell(f_0(X),Y)
  - \ell(f_a(X),Y) \mid Z = z]$. \\

$\widehat\Delta$ &
  Label-free (point) estimator of $\Delta$ or $\Delta(z)$, built without
  target labels. \\

$\varepsilon$ &
  Certificate radius at level $\alpha$: with probability $\geq 1-\alpha$,
  $|\Delta - \widehat\Delta| \leq \varepsilon$ (or $\Delta \geq \widehat\Delta
  - \varepsilon$ for the one-sided form). \\

$\alpha$ &
  Miscoverage level in $(0,1)$; bounds the false-adapt probability
  $\Pr(\textsc{adapt}\ \text{and}\ \Delta \leq 0) \leq \alpha$. \\

$Z$ &
  Label-free observable evidence: $Z = \phi(X_{1:m}, f_0, f_a)$, a vector of
  statistics computable from unlabelled test scores and fixed predictors. \\

$\phi(\cdot)$ &
  Evidence map; extracts $Z$ from the unlabelled batch and predictor pair. \\

$g(Z)$ &
  Decision rule; maps evidence $Z$ to one of $\{\textsc{adapt}, \textsc{freeze},
  \textsc{abstain}\}$. \\

$g^\star$ &
  Bayes-optimal gate: $g^\star(z) = \ind[\Delta(z) > 0]$. \\

$D$ &
  Observable disagreement region: $D = \{x : f_0(x) \neq f_a(x)\}$;
  the support of the sign-of-difference signal. \\

$p_a$ &
  Probability mass of the target distribution on $D$: $p_a = P_T(X \in D)$. \\

$p_0$ &
  Probability mass of the source distribution on $D$: $p_0 = P_S(X \in D)$. \\

$\TV(P, Q)$ &
  Total variation distance between distributions $P$ and $Q$:
  $\TV(P,Q) = \sup_{A} |P(A) - Q(A)|$. \\

$\Law(X \mid P)$ &
  Distribution (law) of $X$ under $P$. \\

$\ind[E]$ &
  Indicator of event $E$; equals $1$ if $E$ holds, $0$ otherwise. \\

$\sign(t)$ &
  Sign function: $\sign(t) = \ind[t > 0]$ (as used in
  Theorem~\ref{thm:gate}). \\

$\argmax$, $\argmin$ &
  Arg-max and arg-min operators. \\

$\R$ &
  The real numbers $(-\infty, +\infty)$. \\

$\E[\cdot]$ &
  Expectation under the distribution clear from context. \\

$\Pr[\cdot]$ &
  Probability under the distribution clear from context. \\

$n$ &
  Sample size (number of paired benefits, calibration tasks, or stream steps
  depending on context). \\

$m$ &
  Size of the unlabelled target batch. \\

$R$ &
  Benefit range $R = b - a$ where the paired benefits $X_i \in [a, b]$. \\

$\widehat{V}$ &
  Unbiased sample variance (denominator $n-1$) of the paired benefits $X_i$. \\

$\widehat\mu$ &
  Empirical mean of the paired benefits; equals $\widehat\Delta$ in the
  batch estimator. \\

$E_t^+$, $E_t^-$ &
  Two one-sided betting e-processes (Theorem~\ref{thm:evalue}):
  $E_t^+$ tests $H_0^+: \Delta \leq 0$; $E_t^-$ tests $H_0^-: \Delta \geq 0$. \\

$\lambda_i$, $\nu_i$ &
  Predictable bet sizes for $E_t^+$ and $E_t^-$ respectively (aGRAPA
  truncated bets, formed from past observations). \\

$\texttt{ks\_mean}$, $\texttt{ks\_max}$ &
  Mean and maximum per-detector Kolmogorov--Smirnov drift between calibration
  and test score marginals (fields of \texttt{Evidence}). \\

$\texttt{disagree}$ &
  $1 - \overline{\rho}$: one minus the mean pairwise rank correlation of
  test scores across detectors (field of \texttt{Evidence}). \\

$\texttt{ess\_frac}$ &
  $\mathrm{ESS}/n_\mathrm{test}$: scale-free effective sample size of
  the Gaussian-ratio importance weights (field of \texttt{Evidence}). \\

$\delta$ &
  Detection threshold or evidence separation; in the quantitative
  non-identifiability theorem (Theorem~\ref{thm:imp}) a location-shift
  parameter indexing $\TV = 2\Phi(\tfrac{1}{2}\sqrt{n}\,\delta) - 1$. \\

Regret-to-oracle &
  $R(g) - R(g^\star)$: excess risk of decision rule $g$ over the
  per-condition oracle $g^\star = \ind[\Delta > 0]$. \\

\end{longtable}

% -----------------------------------------------------------------------
\section{Glossary}
\label{sec:app-notation-glossary}
% -----------------------------------------------------------------------

\begin{description}

\item[Test-time adaptation (TTA).]
  The practice of updating (adapting) a pre-trained model using only unlabelled
  data from a new target distribution at deployment time.  Standard TTA
  mechanisms include entropy minimisation (Tent \cite{wang2021tent}),
  sample-filtered updates (EATA \cite{niu2022eata}), and sharpness-aware resets
  (SAR \cite{niu2023sar}).

\item[Trichotomy.]
  The three-way decision returned by \KGA{}: \textsc{adapt} (adapting is
  certified beneficial), \textsc{freeze} (adapting is certified harmful), or
  \textsc{abstain} (the sign of the benefit is not knowable at level $\alpha$).
  Contrasts with the binary adapt/freeze choice of all prior TTA methods.

\item[Label-free evidence $Z$.]
  The vector of statistics $Z = \phi(X_{1:m}, f_0, f_a)$ computable from
  unlabelled target data, the frozen model, and the candidate adaptation,
  without any target labels.  In the \texttt{kga} package, $Z$ is represented
  by the \texttt{Evidence} dataclass with fields \texttt{ks\_mean},
  \texttt{ks\_max}, \texttt{disagree}, \texttt{entropy\_shift},
  \texttt{conf\_shift}, \texttt{calib\_conf}, \texttt{test\_conf},
  \texttt{ess}, and \texttt{ess\_frac}.

\item[Certificate.]
  A finite-sample statement of the form $\Delta \geq \widehat\Delta - \varepsilon$
  holding with probability at least $1 - \alpha$.  Represented in the \texttt{kga}
  package as the \texttt{Certificate} dataclass with fields \texttt{delta\_hat},
  \texttt{epsilon}, \texttt{method}, \texttt{alpha}, \texttt{n}, and derived
  properties \texttt{lower} / \texttt{upper}.

\item[False adapt.]
  The event that \KGA{} returns \textsc{adapt} but the true benefit is
  non-positive ($\Delta \leq 0$); i.e.\ deploying $f_a$ actually hurts.
  Theorem~\ref{thm:cert} guarantees $\Pr(\text{false adapt}) \leq \alpha$.

\item[False freeze.]
  The symmetric event: \KGA{} returns \textsc{freeze} but $\Delta > 0$ (keeping
  $f_0$ forgoes a genuine benefit).  Also bounded at $\leq \alpha$ by the
  two-sided version of the certificate.

\item[Abstention.]
  The safe default action returned when the certificate brackets zero (i.e.\
  $\widehat\Delta - \varepsilon \leq 0 \leq \widehat\Delta + \varepsilon$).
  \KGA{} keeps the frozen model $f_0$ and logs the evidence for monitoring.
  Abstention is mandatory, not merely conservative, in the non-identifiable
  regime (Theorem~\ref{thm:imp}, Corollary~\ref{sec:t2-cor-abstain}).

\item[Empirical-Bernstein certificate.]
  The batch certificate estimator of Maurer \& Pontil (2009)
  \cite{maurer2009empirical}, the default in \KGA{} (\texttt{method='ebern'}).
  The radius is $\varepsilon = \sqrt{2\widehat{V}\ln(2/\alpha)/n} +
  7R\ln(2/\alpha)/(3(n-1))$, which is strictly tighter than Hoeffding when the
  sample variance $\widehat{V} \ll R^2$.

\item[Conformal certificate (split-conformal).]
  A certificate whose radius is the $(1-\alpha)$-quantile of held-out
  calibration residuals $r_i = |\widehat\Delta_i - \Delta_i|$
  \cite{vovk2005algorithmic,angelopoulos2023conformal}.  The estimator used in
  the main K-Bound cross-task experiments (\texttt{method='conformal'}).

\item[E-value / anytime certificate.]
  The anytime-valid certificate based on two one-sided testing-by-betting
  e-processes \cite{ville1939etude,waudbysmith2024betting}.  Unlike the batch
  and conformal certificates, this certificate is valid simultaneously over all
  stream lengths with no multiplicity correction (\texttt{method='evalue'}).

\item[Covariate shift.]
  The setting where $P_T(X) \neq P_S(X)$ but $P_T(Y \mid X) = P_S(Y \mid X)$.
  Theorem~\ref{thm:pos} identifies this as a provably knowable regime under the
  support-overlap condition: large KS drift combined with adequate importance-weight
  ESS allows \KGA{} to certify the correct action.

\item[Disagreement region.]
  The set $D = \{x : f_0(x) \neq f_a(x)\}$ of inputs on which the frozen and
  adapted predictors differ.  Its probability mass under $P_T$ governs the
  sign-of-difference identifiability (Theorem~\ref{thm:disagree}).

\item[Score archive.]
  The 123-task score cache at \texttt{experiments/elara\_u/score\_archive/},
  the main empirical dataset for K-Bound.  Each task is one \texttt{.npz} file
  with keys \texttt{Sval}, \texttt{yval}, \texttt{Stest}, \texttt{ytest},
  \texttt{det\_names}, \texttt{val\_auc}, and \texttt{domain}.  The archive
  covers ADBench tabular (47 tasks), credit-card fraud, UNSW-NB15 cyber,
  behaviour/HAR, NLP (Enron), time-series (NAB/SMD/TSB-AD), and image benchmarks.

\item[Non-identifiability witness.]
  A pair of target worlds $(P_T^1, P_T^2)$ with $\Law(Z \mid P_T^1) = \Law(Z
  \mid P_T^2)$ and $\sign\Delta^1 = -\sign\Delta^2 \neq 0$, proving that no
  label-free rule can commit correctly in both worlds.  The canonical witness
  (Theorem~\ref{thm:imp}) uses $X \sim \mathcal{N}(0,1)$, $f_0 = \ind[X > 0]$,
  $f_a = \ind[X < 0]$, and flipped labels.

\item[Regret-to-oracle.]
  The excess risk $R(g) - R(g^\star)$ of a decision rule $g$ over the
  per-condition oracle $g^\star = \ind[\Delta > 0]$.  By
  Theorem~\ref{thm:gate}, regret decomposes as
  $\E[|\Delta(Z)| \ind\{\sign\widehat\Delta(Z) \neq \sign\Delta(Z)\}]$;
  the experiments report this quantity as the primary performance metric
  alongside AUROC.

\item[Risk-aligned evidence.]
  Evidence $Z$ is risk-aligned (Definition~3 of Chapter~\ref{ch:formulation})
  if it identifies or brackets the sign of $\Delta(z)$; equivalently, $Z$
  suffices to determine whether adapting helps.  Risk alignment is Assumption~(iii)
  inherited by \KGA{} from Theorem~\ref{thm:cert}.

\item[Mechanism-agnostic wrapper.]
  The design pattern of \KGA{}: the algorithm wraps any existing adaptation
  mechanism (Tent, EATA, SAR, or other) without modifying it.  The
  adaptation mechanism supplies $f_a$; \KGA{} decides whether to deploy it.

\end{description}